\documentclass[10pt,a4paper]{amsart}
\usepackage{amsmath,amssymb,amsthm,graphicx}
\usepackage{verbatim}
\newenvironment{reference}{\comment}{\endcomment}
%\newenvironment{reference}{}{}
\newenvironment{slogan}{\comment}{\endcomment}
\newenvironment{history}{\comment}{\endcomment}
% For commutative diagrams we use Xy-pic
\usepackage[all]{xy}

% We use 2cell for 2-commutative diagrams.
\xyoption{2cell}
% Theorem environments.
%
\theoremstyle{plain}
\newtheorem{theorem}[subsection]{Theorem}
\newtheorem{proposition}[subsection]{Proposition}
\newtheorem{lemma}[subsection]{Lemma}

\theoremstyle{definition}
\newtheorem{definition}[subsection]{Definition}
\newtheorem{example}[subsection]{Example}
\newtheorem{exercise}[subsection]{Exercise}
\newtheorem{situation}[subsection]{Situation}

\theoremstyle{remark}
\newtheorem{remark}[subsection]{Remark}
\newtheorem{remarks}[subsection]{Remarks}

\numberwithin{equation}{subsection}

% Macros
%
\def\lim{\mathop{\mathrm{lim}}\nolimits}
\def\colim{\mathop{\mathrm{colim}}\nolimits}
\def\Spec{\mathop{\mathrm{Spec}}}
\def\Hom{\mathop{\mathrm{Hom}}\nolimits}
\def\Ext{\mathop{\mathrm{Ext}}\nolimits}
\def\SheafHom{\mathop{\mathcal{H}\!\mathit{om}}\nolimits}
\def\SheafExt{\mathop{\mathcal{E}\!\mathit{xt}}\nolimits}
\def\Sch{\mathit{Sch}}
\def\Mor{\mathop{\mathrm{Mor}}\nolimits}
\def\Ob{\mathop{\mathrm{Ob}}\nolimits}
\def\Sh{\mathop{\mathit{Sh}}\nolimits}
\def\NL{\mathop{N\!L}\nolimits}
\def\CH{\mathop{\mathrm{CH}}\nolimits}
\def\proetale{{pro\text{-}\acute{e}tale}}
\def\etale{{\acute{e}tale}}
\def\QCoh{\mathit{QCoh}}
\def\Ker{\mathop{\mathrm{Ker}}}
\def\Im{\mathop{\mathrm{Im}}}
\def\Coker{\mathop{\mathrm{Coker}}}
\def\Coim{\mathop{\mathrm{Coim}}}

% Boxtimes
%
\DeclareMathSymbol{\boxtimes}{\mathbin}{AMSa}{"02}

%
% Macros for moduli stacks/spaces
%
\def\QCohstack{\mathcal{QC}\!\mathit{oh}}
\def\Cohstack{\mathcal{C}\!\mathit{oh}}
\def\Spacesstack{\mathcal{S}\!\mathit{paces}}
\def\Quotfunctor{\mathrm{Quot}}
\def\Hilbfunctor{\mathrm{Hilb}}
\def\Curvesstack{\mathcal{C}\!\mathit{urves}}
\def\Polarizedstack{\mathcal{P}\!\mathit{olarized}}
\def\Complexesstack{\mathcal{C}\!\mathit{omplexes}}
% \Pic is the operator that assigns to X its picard group, usage \Pic(X)
% \Picardstack_{X/B} denotes the Picard stack of X over B
% \Picardfunctor_{X/B} denotes the Picard functor of X over B
\def\Pic{\mathop{\mathrm{Pic}}\nolimits}
\def\Picardstack{\mathcal{P}\!\mathit{ic}}
\def\Picardfunctor{\mathrm{Pic}}
\def\Deformationcategory{\mathcal{D}\!\mathit{ef}}
\usepackage{cleveref}
\crefname{originallemma}{original lemma}{original lemmas}
\Crefname{originallemma}{Original lemma}{Original lemmas}
\makeatletter
\newlabel{lemma-annihilates}{{lemma-annihilates}{source}{}{sourcepointer-1}{}}
\newlabel{lemma-annihilates@cref}{{[originallemma][0][]lemma-annihilates}{[1][0][]source}{}{}{}}
\makeatother
\crefname{originallemma}{original lemma}{original lemmas}
\Crefname{originallemma}{Original lemma}{Original lemmas}
\makeatletter
\newlabel{lemma-compute-tor-Iq}{{lemma-compute-tor-Iq}{source}{}{sourcepointer-2}{}}
\newlabel{lemma-compute-tor-Iq@cref}{{[originallemma][0][]lemma-compute-tor-Iq}{[1][0][]source}{}{}{}}
\makeatother
\crefname{originalreference}{original reference}{original references}
\Crefname{originalreference}{Original reference}{Original references}
\makeatletter
\newlabel{perfect-lemma-spectral-sequence}{{perfect-lemma-spectral-sequence}{source}{}{sourcepointer-3}{}}
\newlabel{perfect-lemma-spectral-sequence@cref}{{[originalreference][0][]perfect-lemma-spectral-sequence}{[1][0][]source}{}{}{}}
\makeatother
\crefname{originalreference}{original reference}{original references}
\Crefname{originalreference}{Original reference}{Original references}
\makeatletter
\newlabel{coherent-lemma-vanishing-nr-affines}{{coherent-lemma-vanishing-nr-affines}{source}{}{sourcepointer-4}{}}
\newlabel{coherent-lemma-vanishing-nr-affines@cref}{{[originalreference][0][]coherent-lemma-vanishing-nr-affines}{[1][0][]source}{}{}{}}
\makeatother
\crefname{originalreference}{original reference}{original references}
\Crefname{originalreference}{Original reference}{Original references}
\makeatletter
\newlabel{coherent-lemma-vanishing-nr-affines-quasi-separated}{{coherent-lemma-vanishing-nr-affines-quasi-separated}{source}{}{sourcepointer-5}{}}
\newlabel{coherent-lemma-vanishing-nr-affines-quasi-separated@cref}{{[originalreference][0][]coherent-lemma-vanishing-nr-affines-quasi-separated}{[1][0][]source}{}{}{}}
\makeatother
\crefname{originalreference}{original reference}{original references}
\Crefname{originalreference}{Original reference}{Original references}
\makeatletter
\newlabel{coherent-lemma-coherent-on-proj}{{coherent-lemma-coherent-on-proj}{source}{}{sourcepointer-6}{}}
\newlabel{coherent-lemma-coherent-on-proj@cref}{{[originalreference][0][]coherent-lemma-coherent-on-proj}{[1][0][]source}{}{}{}}
\makeatother
\crefname{originalreference}{original reference}{original references}
\Crefname{originalreference}{Original reference}{Original references}
\makeatletter
\newlabel{coherent-lemma-recover-tail-graded-module}{{coherent-lemma-recover-tail-graded-module}{source}{}{sourcepointer-7}{}}
\newlabel{coherent-lemma-recover-tail-graded-module@cref}{{[originalreference][0][]coherent-lemma-recover-tail-graded-module}{[1][0][]source}{}{}{}}
\makeatother
\title{A fixed power of the ideal annihilates the Tor modules beyond the number of generators}
\author{The Stacks Project Authors}
\date{stacks-project:local-cohomology.tex:local-snapshot-20261012; GNU FDL 1.2 or later}
\begin{document}
\maketitle
\noindent\textit{Editorial note.} All delivered mathematical text is the author's own, copied from the pinned Stacks Project chapter source sources/stacks-local-cohomology/pinned-original-source.tex; no statement or proof step was rewritten, reordered or removed.\par\medskip
\noindent\textit{Editorial note.} Preamble provenance: the Stacks Project keeps one single shared preamble for every chapter; this card takes it directly from the staged sources/stacks-local-cohomology/preamble.tex as three verbatim spans (lines 8-12, 14-18 and 154-217), so no derived file was written and nothing was added or redefined. The wrapper class and the base packages are supplied by the assembler.\par\medskip
\noindent\textit{Editorial note.} Bibliography provenance: neither the counted statement nor its proof contains a citation (this lane rescans the whole span before writing the card), so the card carries no bibliography and no entry was reconstructed.\par\medskip
\noindent\textit{Editorial note.} Selection provenance: this card was opened under the file-based control contract (CONTROL.md), with the mechanical pre-card gate run on its statement and proof interval before the card was written.\par\medskip
\noindent\textit{Allowed standard prerequisites (editorial).}\begin{itemize}
\item The identification of the derived sections of a twist of the pullback of a quasi-coherent module with the derived tensor product of the module with the corresponding ideal power. Author lemma lemma-compute-tor-Iq, invoked at the start of the counted proof (source line 4682). Turns the Tor modules to be bounded into the cohomology of the twist of the pullback on the blowup.
\item The annihilation of the cohomology sheaves of the pullback in the degrees below the negative of the number of generators by a fixed power of the ideal. Author lemma lemma-annihilates, invoked at the end of the counted proof (source lines 4700-4703). Supplies the single constant that annihilates every term of the filtration.
\end{itemize}\medskip
\ifdefined\hypertarget\hypertarget{sourcepointer-1}{}\fi
\noindent\textit{Source reference.} Original lemma lemma-annihilates: sources/stacks-local-cohomology/pinned-original-source.tex, source lines 4701--4701. Navigation pointer only: the counted text refers to the Stacks item lemma-annihilates, which belongs to another chapter of the project and is not part of this card. The pointer carries the target's own label as its number because this chapter's staged snapshot does not contain the other chapter's tag table; its content is not reproduced here.\par\medskip
\ifdefined\hypertarget\hypertarget{sourcepointer-2}{}\fi
\noindent\textit{Source reference.} Original lemma lemma-compute-tor-Iq: sources/stacks-local-cohomology/pinned-original-source.tex, source lines 4682--4682. Navigation pointer only: the counted text refers to the Stacks item lemma-compute-tor-Iq, which belongs to another chapter of the project and is not part of this card. The pointer carries the target's own label as its number because this chapter's staged snapshot does not contain the other chapter's tag table; its content is not reproduced here.\par\medskip
\ifdefined\hypertarget\hypertarget{sourcepointer-3}{}\fi
\noindent\textit{Source reference.} Original reference perfect-lemma-spectral-sequence: sources/stacks-local-cohomology/pinned-original-source.tex, source lines 4688--4688. Navigation pointer only: the counted text refers to the Stacks item perfect-lemma-spectral-sequence, which belongs to another chapter of the project and is not part of this card. The pointer carries the target's own label as its number because this chapter's staged snapshot does not contain the other chapter's tag table; its content is not reproduced here.\par\medskip
\ifdefined\hypertarget\hypertarget{sourcepointer-4}{}\fi
\noindent\textit{Source reference.} Original reference coherent-lemma-vanishing-nr-affines: sources/stacks-local-cohomology/pinned-original-source.tex, source lines 718--718. Navigation pointer only: the counted text refers to the Stacks item coherent-lemma-vanishing-nr-affines, which belongs to another chapter of the project and is not part of this card. The pointer carries the target's own label as its number because this chapter's staged snapshot does not contain the other chapter's tag table; its content is not reproduced here.\par\medskip
\ifdefined\hypertarget\hypertarget{sourcepointer-5}{}\fi
\noindent\textit{Source reference.} Original reference coherent-lemma-vanishing-nr-affines-quasi-separated: sources/stacks-local-cohomology/pinned-original-source.tex, source lines 4690--4690. Navigation pointer only: the counted text refers to the Stacks item coherent-lemma-vanishing-nr-affines-quasi-separated, which belongs to another chapter of the project and is not part of this card. The pointer carries the target's own label as its number because this chapter's staged snapshot does not contain the other chapter's tag table; its content is not reproduced here.\par\medskip
\ifdefined\hypertarget\hypertarget{sourcepointer-6}{}\fi
\noindent\textit{Source reference.} Original reference coherent-lemma-coherent-on-proj: sources/stacks-local-cohomology/pinned-original-source.tex, source lines 4612--4612. Navigation pointer only: the counted text refers to the Stacks item coherent-lemma-coherent-on-proj, which belongs to another chapter of the project and is not part of this card. The pointer carries the target's own label as its number because this chapter's staged snapshot does not contain the other chapter's tag table; its content is not reproduced here.\par\medskip
\ifdefined\hypertarget\hypertarget{sourcepointer-7}{}\fi
\noindent\textit{Source reference.} Original reference coherent-lemma-recover-tail-graded-module: sources/stacks-local-cohomology/pinned-original-source.tex, source lines 4613--4613. Navigation pointer only: the counted text refers to the Stacks item coherent-lemma-recover-tail-graded-module, which belongs to another chapter of the project and is not part of this card. The pointer carries the target's own label as its number because this chapter's staged snapshot does not contain the other chapter's tag table; its content is not reproduced here.\par\medskip
\begin{lemma}
\label{lemma-annihilates-tors}
In the situation above, let $t$ be an upper bound on the number of
generators for $I$. There exists an integer $c = c(A, I) \geq 0$
such that for any $A$-module $M$ the tor modules
$\text{Tor}_i^A(M, A/I^q)$ are annihilated by $I^c$ for $i > t$
and all $q \geq 0$.
\end{lemma}

\noindent
For the rest of the discussion in this section we fix a Noetherian ring $A$
and an ideal $I \subset A$. We denote
$$
p : X \to \Spec(A)
$$
the blowing up of $\Spec(A)$ in the ideal $I$. In other words, $X$ is the
$\text{Proj}$ of the Rees algebra $\bigoplus_{n \geq 0} I^n$.
By Cohomology of Schemes, Lemmas \ref{coherent-lemma-coherent-on-proj} and
\ref{coherent-lemma-recover-tail-graded-module}
we can choose an integer $q(A, I) \geq 0$ such that for all $q \geq q(A, I)$
we have $H^i(X, \mathcal{O}_X(q)) = 0$ for
$i > 0$ and $H^0(X, \mathcal{O}_X(q)) = I^q$.

\begin{proof}
Let $q(A, I)$ be as above. For $q \geq q(A, I)$ we have
$$
R\Gamma(X, Lp^*\widetilde{M}(q)) = M \otimes_A^\mathbf{L} I^q
$$
by Lemma \ref{lemma-compute-tor-Iq}.
We have a bounded and convergent spectral sequence
$$
H^a(X, H^b(Lp^*\widetilde{M}(q))) \Rightarrow
\text{Tor}_{-a - b}^A(M, I^q)
$$
by Derived Categories of Schemes, Lemma \ref{perfect-lemma-spectral-sequence}.
Let $d$ be an integer as in Cohomology of Schemes, Lemma
\ref{coherent-lemma-vanishing-nr-affines-quasi-separated}
(actually we can take $d = t$, see
Cohomology of Schemes, Lemma \ref{coherent-lemma-vanishing-nr-affines}).
Then we see that $H^{-i}(X, Lp^*\widetilde{M}(q)) = \text{Tor}_i^A(M, I^q)$
has a finite filtration with at most $d$ steps whose graded are
subquotients of the modules
$$
H^a(X, H^{- i - a}(Lp^*\widetilde{M})(q)),\quad
a = 0, 1, \ldots, d - 1
$$
If $i \geq t$ then all of these modules are annihilated
by $I^c$ where $c = c(A, I)$ is as in Lemma \ref{lemma-annihilates}
because the cohomology sheaves $H^{- i - a}(Lp^*\widetilde{M})$
are all annihilated by $I^c$ by the lemma. Hence we see that
$\text{Tor}_i^A(M, I^q)$ is annihilated by $I^{dc}$ for
$q \geq q(A, I)$ and $i \geq t$. Using the short exact sequence
$0 \to I^q \to A \to A/I^q \to 0$ we find that
$\text{Tor}_i(M, A/I^q)$ is annihilated by $I^{dc}$ for $q \geq q(A, I)$
and $i > t$. We conclude that $I^m$ with $m = \max(dc, q(A, I) - 1)$
annihilates $\text{Tor}_i^A(M, A/I^q)$ for all $q \geq 0$
and $i > t$ as desired.
\end{proof}

\end{document}
